% Generated draft: requires mathematical and source review.
\documentclass[11pt]{article}
\usepackage[margin=25mm]{geometry}
\usepackage{fontspec}
\setmainfont{Latin Modern Roman}
\usepackage{amsmath,amssymb,mathtools,unicode-math}
\setmathfont{Latin Modern Math}
\usepackage{enumitem,needspace,xurl,hyperref}
\setlength{\parindent}{0pt}
\setlength{\parskip}{5pt}
\setlength{\emergencystretch}{3em}
\newcommand{\N}{\mathbb{N}}
\newcommand{\Z}{\mathbb{Z}}
\newcommand{\Q}{\mathbb{Q}}
\newcommand{\R}{\mathbb{R}}
\newcommand{\C}{\mathbb{C}}
\newcommand{\F}{\mathbb{F}}
\newcommand{\A}{\mathbb{A}}
\newcommand{\PP}{\mathbb{P}}
\newcommand{\E}{\mathbb{E}}
\newcommand{\eps}{\varepsilon}
\newcommand{\dd}{\,\mathrm{d}}
\begin{document}
\section*{UnsolvedMath ID: 30005785}
\textbf{Source problem code:} catalogue entry
\subsection{Definitions and mathematical statement}
Let \(\Gamma\) be a finite simplicial graph with vertex set \(\{1,\ldots,n\}\), and assign an integer \(m_{ij}\geq 2\) to each edge \(\{i,j\}\). Set \(m_{ij}=\infty\) when distinct vertices \(i,j\) are not joined by an edge. For letters \(a,b\) and a positive integer \(m\), let \(\operatorname{Alt}_m(a,b)\) denote the alternating word of length \(m\) beginning with \(a\). The associated Artin group is
\[
A_\Gamma=\left\langle s_1,\ldots,s_n\ \middle|\ \operatorname{Alt}_{m_{ij}}(s_i,s_j)=\operatorname{Alt}_{m_{ij}}(s_j,s_i)\text{ whenever }i<j\text{ and }m_{ij}<\infty\right\rangle.
\]
There is no defining relation for a pair with \(m_{ij}=\infty\).

The group \(A_\Gamma\) is torsion-free if
\[
\forall g\in A_\Gamma\ \forall k\in\Z_{\geq 1},\qquad g^k=1\Longrightarrow g=1.
\]
It has a solvable word problem if there exists an algorithm which, on every finite word in \(s_1^{\pm1},\ldots,s_n^{\pm1}\), terminates and decides whether that word represents the identity in \(A_\Gamma\). This is an assertion for each fixed \(\Gamma\); it does not additionally require a single algorithm taking arbitrary labeled graphs as input.

To specify the conventional meaning of the third property, let
\[
W_\Gamma=A_\Gamma/\langle\!\langle s_i^2\mid 1\leq i\leq n\rangle\!\rangle
\]
be the associated Coxeter group. A subset \(T\subseteq\{1,\ldots,n\}\) is spherical if the subgroup of \(W_\Gamma\) generated by the images of \(s_i\), \(i\in T\), is finite. Let \(S_\Gamma\) be the standard Salvetti complex associated with this Coxeter system: the usual connected CW complex with one cell of dimension \(|T|\) for each spherical subset \(T\), with the standard Coxeter--Salvetti attaching maps. Its fundamental group is naturally \(A_\Gamma\). The Artin-group \(K(\pi,1)\) conjecture for \(\Gamma\) asserts that \(S_\Gamma\) is a \(K(A_\Gamma,1)\), equivalently that its universal cover is contractible, or that
\[
\pi_q(S_\Gamma)=0\qquad\text{for every integer }q\geq 2.
\]
Here the specified standard complex matters: merely asking whether some \(K(A_\Gamma,1)\) exists would not express this conjecture.

The supplied reconstructed challenge asks whether, for every finite labeled graph \(\Gamma\) as above, all three properties hold. Equivalently, it asks whether each of the following collections of labeled graphs is empty:
\[
\begin{aligned}
\mathcal F_{\mathrm{tor}}&=\{\Gamma:A_\Gamma\text{ has a nonidentity element of finite order}\},\\
\mathcal F_{\mathrm{word}}&=\{\Gamma:A_\Gamma\text{ has an unsolvable word problem}\},\\
\mathcal F_{K}&=\{\Gamma:S_\Gamma\text{ is not a }K(A_\Gamma,1)\}.
\end{aligned}
\]
If any collection is nonempty, the alternative requested in the supplied statement is a precise characterization of that collection, separately for each property.

This formulation records the supplied reconstructed statement. The supplied original excerpt contains only a general observation about long-standing questions and a damaged Artin presentation; it does not explicitly state these three conjectures or specify the Salvetti formulation. The dated supplied assessment likewise describes the entry as a general warning rather than a single theorem. Thus the intended three-part target and its correspondence with the report require human source review.
\subsection{Short English statement}
For every Artin group defined by a finite labeled graph, the supplied reconstructed challenge asks for torsion-freeness, an algorithm deciding whether a word represents the identity, and asphericity of the standard Salvetti complex. If any property fails, it asks for an exact characterization of the graphs for which it fails. The supplied original excerpt does not explicitly identify these three targets, so this remains a statement-only interpretation requiring review.
\subsection{Sources}
\begin{itemize}
\item Supplied local UnsolvedMath registry record, permanent ID \(30005785\), title: \emph{Three Structural Problems for Artin Groups}; collection: \emph{Oberwolfach Reports Open Problems}; exact source problem code: \texttt{catalogue entry}. The record supplies proposed year 2024 and no named proposer.
\item Supplied source citation: \emph{Mini-Worskhop: Artin Groups meet Triangulated Categories} (2024), Oberwolfach Reports, DOI 10.4171/owr/2024/4, pp. 203--234. Supplied source URL: \url{https://doi.org/10.4171/owr/2024/4}. The damaged original excerpt includes a page marker for p. 217.
\item Supplied literature source, labeled ``original OWR report'': \url{https://doi.org/10.4171/owr/2024/4}.
\item Literature assessment dated exactly \texttt{2026-08-22}: the supplied curation labels the entry open and describes it as a general warning rather than a single theorem, concerning several long-standing Artin-group problems remaining open outside special classes. This is supplied dated curation, not an independently verified literature assessment; its time-sensitive status and cited claims require independent checking.
\item The registry labels the statement \texttt{corrected\_verified} and reports that it was reconstructed and checked against the official Oberwolfach report. These are supplied provenance claims, not verification performed for this draft.
\end{itemize}
\end{document}
